\section{Bases of normal tensors}

\begin{definition}[Basis of Normal Tensors (BNT)]\label{def:bnt}
    \lean{MPSTensor.IsBNT}
    \leanok
    \uses{def:normal, def:mpv}
    A \emph{basis of normal tensors} (BNT) for a total tensor $A_{\mathrm{tot}}$
    consists of $g$ block tensors $(A_1, \ldots, A_g)$ with bond dimensions
    $(D_1, \ldots, D_g)$ such that each $A_j$ is normal
    (Definition~\ref{def:normal}); at each system size~$N\geq1$ the MPV of
    $A_{\mathrm{tot}}$ lies in the span of the block MPVs
    $\{\ket{V^{(N)}(A_j)}\}_{j=1}^g$; and for sufficiently large~$N$ those block
    MPVs are linearly independent. This is \cite[Definition~IV.2]{Cirac2021Matrix}.
    In the canonical-form characterization of
    \cite[Section~2.3]{Cirac2016MPDO_arXiv}, a BNT is a minimal family of
    representatives for $\widetilde A_k^i = e^{i\phi_k} X_k A_j^i X_k^{-1}$ among
    the normal blocks.
\end{definition}

\begin{definition}[Basis of normal tensors, coefficient form]\label{def:bnt_cpsv}
    \lean{MPSTensor.IsCPSVBasisOfNormalTensors}
    \leanok
    \uses{def:nt_cpsv, def:mpv}
    % Source: Cirac2017MPDO_AnnPhys \S~II.C, lines 271--274.
    The coefficient form of a basis of normal tensors: each block $A_j$ is a
    CPSV normal tensor (Definition~\ref{def:nt_cpsv}), the MPV of $A$ is given at
    every positive length by explicit coefficients $c_j^{(N)} \in \C$ with
    $\ket{V^{(N)}(A)} = \sum_{j=1}^{g} c_j^{(N)} \ket{V^{(N)}(A_j)}$, and the
    block MPVs $\{\ket{V^{(N)}(A_j)}\}$ are eventually linearly independent.
\end{definition}

\begin{lemma}[Positive bond dimensions in a coefficient-form BNT]
    \label{lem:cpsv_bnt_block_dimensions_positive}
    \lean{MPSTensor.IsCPSVBasisOfNormalTensors.blocks_dim_pos}
    \leanok
    \uses{def:bnt_cpsv}
    Every block in a coefficient-form basis of normal tensors has positive
    bond dimension.
\end{lemma}

\begin{proof}\leanok
    \uses{def:bnt_cpsv}
    Eventual linear independence makes the matrix product vector of each block
    nonzero at some positive length.  A tensor with zero bond dimension has
    zero matrix product vector at every length, which is a contradiction.
\end{proof}

\begin{lemma}[Bond dimension one implies irreducibility]
    \label{lem:bond_dimension_one_irreducible}
    \lean{MPSTensor.isIrreducibleTensor_of_bondDim_one}
    \leanok
    \uses{def:irreducible_tensor}
    Every tensor with a one-dimensional bond space is irreducible.
\end{lemma}

\begin{proof}\leanok
    The only orthogonal projections on a one-dimensional complex vector space
    are zero and the identity.
\end{proof}

\begin{lemma}[Normality in bond dimension one]
    \label{lem:bond_dimension_one_identity_transfer_normal}
    \lean{MPSTensor.isNormalTensor_of_bondDim_one_of_transferMap_eq_id}
    \leanok
    \uses{def:nt_cpsv, lem:bond_dimension_one_irreducible}
    A tensor with one-dimensional bond space and identity transfer map is a
    normal tensor.
\end{lemma}

\begin{proof}\leanok
    Irreducibility follows from the preceding lemma. The identity transfer map
    has spectral radius one, and its only eigenvalue is one.
\end{proof}

\begin{theorem}[Distinct BNT members are not gauge-phase equivalent]
    \label{thm:cpsv_bnt_blocks_not_gaugePhaseEquiv}
    \lean{MPSTensor.IsCPSVBasisOfNormalTensors.blocks_not_gaugePhaseEquiv}
    \leanok
    \uses{def:bnt_cpsv, def:blocks_not_gauge_phase_equiv}
    Distinct members of a basis of normal tensors are not gauge-phase
    equivalent after identifying equal bond dimensions.
\end{theorem}

\begin{proof}\leanok
    At every sufficiently large length the matrix product vectors of the basis
    members are linearly independent. A gauge-phase equivalence between two
    distinct members would give, for a nonzero scalar $\zeta$,
    \[
        \ket{V^{(N)}(A_k)}
        = \zeta^N\ket{V^{(N)}(A_j)}.
    \]
    Thus one vector would be a scalar multiple of the other, contradicting
    linear independence.
\end{proof}

\begin{theorem}[Characterization by canonical-form representatives]
    \label{thm:cpsv_bnt_characterization}
    \lean{MPSTensor.isCPSVBasisOfNormalTensors_iff_canonicalForm_covered_and_minimal}
    \leanok
    \uses{def:bnt_cpsv, def:nt_cpsv, def:block_diagonal_tensor,
        def:same_mpv2_pos, def:blocks_not_gauge_phase_equiv}
    Let $A$ have a canonical-form MPV expansion
    \[
        \ket{V^{(N)}(A)}
        = \sum_{k=1}^{r}\mu_k^N\ket{V^{(N)}(\widetilde A_k)}
    \]
    at every positive length, where the tensors $\widetilde A_k$ are normal
    tensors of positive bond dimension. Set
    $K_{\mathrm{act}}=\{k:\mu_k\ne0\}$. Let $(A_j)_{j=1}^{g}$ be tensors of
    positive bond dimension. Then
    $(A_j)_{j=1}^{g}$ is a basis of normal tensors for $A$ if and only if:
    \begin{enumerate}
        \item every $A_j$ is normal;
        \item for every $k\in K_{\mathrm{act}}$ there are an index $j$, an
              invertible matrix $X_k$, and a phase $e^{i\phi_k}$ such that
              \[
                  \widetilde A_k^i
                  = e^{i\phi_k}X_kA_j^iX_k^{-1}
              \]
              for every physical index $i$;
        \item no two distinct tensors $A_j$ and $A_{j'}$ are related in this
              way.
    \end{enumerate}
    This is \cite[Proposition~2.7]{Cirac2016MPDO_arXiv}, with the literal
    canonical-form display restricted to its active blocks.
    % Local correction documented in
    % docs/paper-gaps/cpsv16_bnt_characterization_active_blocks.tex.
\end{theorem}

\begin{proof}\leanok
    \uses{thm:cpsv_normal_mpv_phase_gauge,
        thm:cpsv_normal_mpv_phase_alternative,
        thm:cpsv_bnt_blocks_not_gaugePhaseEquiv,
        lem:bnt_finite_index_overlap_orthonormal_li,
        lem:finite_geometric_sum_eventual_vanishing,
        lem:sector_decomposition_coeff_not_eventually_zero,
        def:mpv_phase_class_data,
        thm:cpsv_normal_separated_eventually_li,
        thm:mpv_decomposition}
    First discard the zero-weight summands; they contribute zero at every
    positive length. Reindex the remaining summands, group them into the MPV
    phase classes of Definition~\ref{def:mpv_phase_class_data}, and choose one
    representative from each class. By
    Theorem~\ref{thm:cpsv_bnt_blocks_not_gaugePhaseEquiv}, eventual linear
    independence of a BNT excludes gauge-phase equivalence between two
    distinct $A_j$.

    Suppose that some phase-class representative is not gauge-phase equivalent
    to any $A_j$. Partition all phase classes into the unmatched classes $T$
    and the remaining classes. Every remaining representative is a nonzero
    scalar multiple of some $\ket{V^{(N)}(A_j)}$ at each positive length.
    Distinct unmatched representatives have vanishing mutual overlaps, and
    their overlaps with every $A_j$ also vanish. Explicitly, for $t\ne t'$ in
    $T$ and every $j$,
    \[
        O_{\widetilde A_t\widetilde A_{t'}}(N)\longrightarrow0,
        \qquad
        O_{\widetilde A_t A_j}(N)\longrightarrow0.
    \]
    By Lemma~\ref{lem:bnt_finite_index_overlap_orthonormal_li}, the combined
    family
    \[
        \{\ket{V^{(N)}(\widetilde A_t)}:t\in T\}
        \cup \{\ket{V^{(N)}(A_j)}:1\leq j\leq g\}
    \]
    is linearly independent for all sufficiently large $N$. Comparing the
    canonical-form expansion with the BNT expansion gives a relation
    \[
        \sum_{t\in T}c_{N,t}\ket{V^{(N)}(\widetilde A_t)}
        +\sum_{j=1}^{g}\alpha_{N,j}\ket{V^{(N)}(A_j)}=0.
    \]
    Linear independence forces every $c_{N,t}$ to vanish. Hence the
    coefficient of every unmatched phase class vanishes for all sufficiently
    large $N$.
    For an unmatched class $t$, let $q$ range over its active canonical-form
    summands, write $\mu_q$ for the corresponding canonical-form weight, and
    write $\zeta_q$ for the phase relating that summand to the chosen
    representative. Then
    \[
        c_{N,t}=\sum_{q\in t}(\zeta_q\mu_q)^N.
    \]
    Every base $\zeta_q\mu_q$ is nonzero because only active summands occur.
    Eventual vanishing of this finite geometric sum would imply vanishing at
    $N=0$, where its value is the positive number of summands in the class.
    This contradiction proves coverage. The normal-tensor phase theorem turns
    the resulting MPV phase relation into the stated gauge-phase relation.

    Conversely, coverage gives, for each canonical-form summand,
    \[
        \ket{V^{(N)}(\widetilde A_k)}
        = \zeta_k^N\ket{V^{(N)}(A_{j(k)})}.
    \]
    Hence
    \[
        \ket{V^{(N)}(A)}
        = \sum_j
            \left(\sum_{k:j(k)=j}(\mu_k\zeta_k)^N\right)
            \ket{V^{(N)}(A_j)},
    \]
    which is the required BNT span at every positive length. Pairwise
    gauge-phase distinction and normality give eventual linear independence,
    completing the converse.
\end{proof}

\begin{theorem}[Failure of literal BNT uniqueness]
    \label{thm:cpsv_bnt_uniqueness_counterexample}
    \lean{MPSTensor.exists_isCPSVBasisOfNormalTensors_unequal_card}
    \leanok
    \uses{def:bnt_cpsv, def:nt_cpsv, def:same_mpv2_pos,
        def:block_diagonal_tensor}
    There is a normal tensor with a one-block canonical-form expansion which
    admits two bases of normal tensors of different cardinalities. Hence the
    unrestricted uniqueness sentence of
    \cite[Appendix~A, line~1148]{Cirac2016MPDO_arXiv} does not follow from the
    literal definition of a basis of normal tensors.
\end{theorem}

\begin{proof}\leanok
    \uses{thm:mpv_decomposition}
    Take physical dimension two and bond dimension one. Let $A_0$ be supported
    on physical symbol zero and let $A_1$ be supported on physical symbol one:
    \[
        A_0^0=1,\quad A_0^1=0,
        \qquad
        A_1^0=0,\quad A_1^1=1.
    \]
    Both transfer maps are the identity on the one-dimensional matrix algebra:
    \[
        \mathcal E_{A_j}(X)=X \qquad (j=0,1).
    \]
    Hence both tensors are normal. At every positive length,
    \[
        \ket{V^{(N)}(A_0)}=\ket{0\cdots0},
        \qquad
        \ket{V^{(N)}(A_1)}=\ket{1\cdots1}.
    \]
    Thus $(A_0)$ and $(A_0,A_1)$ are both bases of normal tensors for $A_0$.
    Indeed, for every $N\geq 1$, the second expansion is
    \[
        \ket{V^{(N)}(A_0)}
        =1\cdot\ket{V^{(N)}(A_0)}+0\cdot\ket{V^{(N)}(A_1)}.
    \]
    Their index sets have cardinalities one and two, so they cannot be related
    by a bijection.
    % The missing converse-coverage condition is documented in
    % docs/paper-gaps/cpsv16_bnt_uniqueness_zero_coefficient.tex.
\end{proof}

\begin{theorem}[Restricted BNT uniqueness under converse coverage]
    \label{thm:cpsv_bnt_uniqueness_converse_coverage}
    \lean{MPSTensor.IsCPSVBasisOfNormalTensors.equiv_of_converse_coverage}
    \leanok
    \uses{def:bnt_cpsv, def:nt_cpsv, def:same_mpv2_pos,
        def:block_diagonal_tensor, def:gauge_phase_equiv}
    Let $A$ have a canonical-form MPV expansion
    \[
        \ket{V^{(N)}(A)}
        =\sum_{k=0}^{r-1}\mu_k^N\ket{V^{(N)}(\widetilde A_k)}
    \]
    at every positive length, where the tensors $\widetilde A_k$ are normal
    and have positive bond dimensions.
    Let $(A_j)_{j=0}^{g_1-1}$ and $(B_l)_{l=0}^{g_2-1}$ be two bases of normal
    tensors for $A$, all of whose members have positive bond dimensions.
    Assume that every $A_j$ and every $B_l$ is related by a phase and an
    invertible gauge to a block $\widetilde A_k$ with $\mu_k\ne0$. Then there
    is a bijection
    $e:\{0,\ldots,g_1-1\}\to\{0,\ldots,g_2-1\}$ such that the bond dimensions of
    $A_j$ and $B_{e(j)}$ agree and
    \[
        B_{e(j)}^i=\zeta_jX_jA_j^iX_j^{-1},
        \qquad |\zeta_j|=1,
    \]
    for every $j$ and every physical index $i$.
\end{theorem}

\begin{proof}\leanok
    \uses{thm:cpsv_bnt_characterization,
        thm:cpsv_bnt_blocks_not_gaugePhaseEquiv,
        thm:cpsv_normal_mpv_phase_gauge,
        lem:gauge_phase_equiv_cast_compose, def:gauge_phase_equiv}
    For each $A_j$, choose an active canonical-form block
    $\widetilde A_{\ell(j)}$ representing it, and then choose a member
    $B_{f(j)}$ representing that block. If $f(j)=f(j')$, then
    \[
        A_j\sim_{\mathrm{gp}}\widetilde A_{\ell(j)}
        \sim_{\mathrm{gp}}B_{f(j)}
        =B_{f(j')}\sim_{\mathrm{gp}}\widetilde A_{\ell(j')}
        \sim_{\mathrm{gp}}A_{j'}.
    \]
    Minimality of the first basis gives $j=j'$, so $f$ is injective.
    Reversing the roles of the two bases gives an injection in the other
    direction. Since both index sets are finite, $f$ is bijective; set $e=f$.
    Write the two relations through the chosen active block as
    \[
        \widetilde A_{\ell(j)}^i
        =\zeta_{1,j}X_{1,j}A_j^iX_{1,j}^{-1},
        \qquad
        \widetilde A_{\ell(j)}^i
        =\zeta_{2,j}X_{2,j}B_{e(j)}^iX_{2,j}^{-1}.
    \]
    Solving the second identity for $B_{e(j)}^i$ and using the first gives
    \[
        B_{e(j)}^i
        =(\zeta_{2,j}^{-1}\zeta_{1,j})
        (X_{2,j}^{-1}X_{1,j})A_j^i(X_{1,j}^{-1}X_{2,j}).
    \]
    Denote the scalar and gauge in this identity by $\zeta_j$ and $X_j$.
    The gauge-phase relation gives
    \[
        O_{B_{e(j)}B_{e(j)}}(N)
        =(\zeta_j\overline{\zeta_j})^N O_{A_jA_j}(N)
        =|\zeta_j|^{2N}O_{A_jA_j}(N).
    \]
    Normality gives
    $O_{A_jA_j}(N)\to1$ and $O_{B_{e(j)}B_{e(j)}}(N)\to1$; hence the last
    identity forces $|\zeta_j|=1$.
    % This is the converse-coverage restriction documented in
    % docs/paper-gaps/cpsv16_bnt_uniqueness_zero_coefficient.tex, not the
    % unrestricted sentence at arXiv:1606.00608, Appendix A, line 1148.
\end{proof}

\begin{lemma}[Perron gauges preserve primitivity]
    \label{lem:perron_gauge_primitive_iff}
    \lean{MPSTensor.isPrimitive_transferMap_tpGauge_iff}
    \leanok
    \uses{def:primitive_channel}
    Let $\sigma$ be positive definite. The transfer map of the Perron gauge
    determined by $\sigma$ is primitive if and only if the original transfer
    map is primitive.
\end{lemma}

\begin{proof}\leanok
    The two transfer maps are similar through the invertible congruence
    $X\mapsto\sigma^{-1/2}X\sigma^{-1/2}$. Similar linear maps have the same
    spectrum, so the condition that $1$ is the only eigenvalue on the unit
    circle is unchanged.
\end{proof}

\begin{theorem}[Spectral normality gives a pure trace-preserving gauge]
    \label{thm:cpsv_normal_tp_gauge}
    \lean{MPSTensor.IsNormalTensor.exists_tpGauge}
    \leanok
    \uses{def:nt_cpsv}
    Let $A$ be a normal tensor of positive bond dimension in the sense of
    \cite[Section~2.2]{Cirac2016MPDO_arXiv}. There is a positive definite
    matrix $\sigma$ such that the gauge obtained from $\sigma$ is
    trace-preserving, primitive, and irreducible. No scalar rescaling is
    needed.
\end{theorem}

\begin{proof}\leanok
    \uses{lem:perron_gauge_primitive_iff}
    The irreducible Perron--Frobenius theorem gives positive definite right and
    left eigenvectors with positive eigenvalues $r$ and $t$. The spectral
    normalization gives $r=1$. Adjointness with respect to the trace pairing
    gives
    \[
        r\,\Tr(\sigma\rho)=t\,\Tr(\sigma\rho).
    \]
    Since $\sigma$ and $\rho$ are positive definite, their pairing is nonzero,
    hence $t=1$. Thus the Perron gauge is a pure similarity. Irreducibility and
    primitivity are invariant under this similarity.
\end{proof}

\begin{theorem}[Spectral normality gives eventual block injectivity]
    \label{thm:cpsv_normal_eventually_injective}
    \lean{MPSTensor.IsNormalTensor.isNormal}
    \leanok
    \uses{def:nt_cpsv, def:normal}
    Every positive-bond-dimension normal tensor in the sense of
    \cite[Section~2.2]{Cirac2016MPDO_arXiv} is injective after blocking by some
    finite length.
\end{theorem}

\begin{proof}\leanok
    \uses{thm:cpsv_normal_tp_gauge, thm:normal_tp_primitive_irred}
    Apply Theorem~\ref{thm:cpsv_normal_tp_gauge}. The trace-preserving,
    primitive, irreducible representative is eventually block-injective, and
    gauge invariance transports this conclusion back to the original tensor.
\end{proof}

\begin{theorem}[A coefficient-form BNT is an algebraic BNT]
    \label{thm:cpsv_bnt_is_bnt}
    \lean{MPSTensor.IsCPSVBasisOfNormalTensors.isBNT}
    \leanok
    \uses{def:bnt, def:bnt_cpsv,
        lem:cpsv_bnt_block_dimensions_positive,
        thm:cpsv_normal_eventually_injective}
    If $(A_j)_{j=1}^{g}$ is a coefficient-form basis of normal tensors for
    $A$, then $(A_j)_{j=1}^{g}$ is an algebraic basis of normal tensors for
    $A$. Only this implication from spectral normality to algebraic normality
    is asserted.
\end{theorem}

\begin{proof}\leanok
    \uses{lem:cpsv_bnt_block_dimensions_positive,
        thm:cpsv_normal_eventually_injective}
    Every block has positive bond dimension, so spectral normality implies
    eventual block injectivity for each block. The spanning identities and
    eventual linear independence are unchanged.
\end{proof}

\begin{lemma}[Transport of an algebraic BNT along positive-length MPV equality]
    \label{lem:bnt_of_same_mpv2_pos}
    \lean{MPSTensor.IsBNT.of_sameMPV₂Pos}
    \leanok
    \uses{def:bnt, def:same_mpv2_pos}
    Let $A$ and $A'$ have the same matrix product vectors at every positive
    length. If $(A_j)_{j=1}^{g}$ is a basis of normal tensors for $A$, then
    the same family is a basis of normal tensors for $A'$.
\end{lemma}

\begin{proof}\leanok
    \uses{def:bnt, def:same_mpv2_pos}
    Substitute the positive-length MPV equality into the spanning identity
    for $A$. Normality and eventual linear independence depend only on the
    family $(A_j)_{j=1}^{g}$.
\end{proof}

\begin{theorem}[Non-decaying rectangular overlap determines the bond dimension]
    \label{thm:rect_overlap_norm_one_dim_eq}
    \lean{MPSTensor.dim_eq_of_mpvOverlap_norm_tendsto_one_of_irreducible_TP}
    \leanok
    \uses{def:mpv_overlap, def:irreducible_tensor, def:tp_kraus}
    Let $A$ and $B$ be irreducible trace-preserving tensors of positive,
    possibly different, bond dimensions. If
    \[
        |O_{AB}(N)|\longrightarrow 1,
    \]
    then their bond dimensions are equal.
\end{theorem}

\begin{proof}\leanok
    \uses{thm:overlap_decay_rect_irr_tp}
    If the bond dimensions differed, the rectangular overlap-decay theorem
    would give $O_{AB}(N)\to0$, and hence $|O_{AB}(N)|\to0$, contradicting the
    assumed limit.
\end{proof}

\begin{theorem}[Phase-equivalent normal tensors are gauge-phase equivalent]
    \label{thm:cpsv_normal_mpv_phase_gauge}
    \lean{MPSTensor.IsNormalTensor.exists_apply_ne_zero}
    \lean{MPSTensor.IsNormalTensor.selfOverlap_tendsto_one}
    \lean{MPSTensor.isNormalTensor_cast_iff}
    \lean{MPSTensor.norm_eq_one_of_gaugePhase_cast_of_isNormalTensor}
    \lean{MPSTensor.MPVBlockPhaseEquiv.dim_eq_and_gaugePhaseEquiv_of_isNormalTensor}
    \leanok
    \uses{def:nt_cpsv, def:mpv_block_phase_equiv,
        def:gauge_phase_equiv}
    Let $A$ and $B$ be normal tensors of positive, possibly different, bond
    dimensions. Suppose that, for some nonzero $\zeta\in\C$,
    \[
        \ket{V^{(N)}(B)}=\zeta^N\ket{V^{(N)}(A)}
        \qquad (N\geq 0).
    \]
    Then the bond dimensions are equal and there is an invertible matrix $X$
    and a scalar $\eta$ of unit modulus such that
    \[
        B^i=\eta X A^iX^{-1}
    \]
    for every physical letter $i$. This is the gauge recovery used in
    \cite[Appendix~A]{Cirac2016MPDO_arXiv}.
\end{theorem}

\begin{proof}\leanok
    \uses{thm:cpsv_normal_tp_gauge, thm:rect_overlap_norm_one_dim_eq,
        thm:nt_overlap_norm_one_gauge}
    Apply Theorem~\ref{thm:cpsv_normal_tp_gauge} to both tensors. Gauge
    invariance of matrix product vectors preserves the phase relation, whose
    self-overlap identities imply unit modulus for the phase. The resulting
    rectangular cross-overlap has norm tending to one. If the bond dimensions
    differed, Theorem~\ref{thm:rect_overlap_norm_one_dim_eq} would give a
    contradiction. Thus the bond dimensions agree, and the
    unit-overlap rigidity theorem gives a gauge-phase equivalence between the
    trace-preserving representatives. Conjugating by the two Perron gauges
    gives the asserted relation for $A$ and $B$.
\end{proof}

\begin{theorem}[Normal-tensor overlap dichotomy]
    \label{thm:cpsv_normal_overlap_dichotomy}
    \lean{MPSTensor.IsNormalTensor.overlap_dichotomy}
    \leanok
    \uses{def:nt_cpsv, def:mpv_overlap, def:gauge_phase_equiv}
    Let $A$ and $B$ be normal tensors of positive, possibly different, bond
    dimensions. Their self-overlaps satisfy
    \[
        O_{AA}(N)\longrightarrow 1,
        \qquad O_{BB}(N)\longrightarrow 1.
    \]
    Moreover, either
    \[
        O_{AB}(N)\longrightarrow 0,
    \]
    or their bond dimensions are equal; call the common value $D$. Then there
    are $X\in\GL_D(\C)$ and $\zeta\in\C$, with $|\zeta|=1$, such that
    \[
        B^i=\zeta X A^iX^{-1}
    \]
    for every physical letter $i$. In the latter case,
    $|O_{AB}(N)|\to 1$. This is
    \cite[Lemma~A.2]{Cirac2016MPDO_arXiv}.
\end{theorem}

\begin{proof}\leanok
    \uses{thm:cpsv_normal_tp_gauge, thm:overlap_decay_rect_irr_tp,
        thm:overlap_decay_irr_tp, thm:gauge_invariance,
        thm:gauge_phase_mpv_scaling, thm:cpsv_normal_mpv_phase_gauge}
    Put both tensors in their pure trace-preserving gauges. If their bond
    dimensions differ, rectangular overlap decay gives $O_{AB}(N)\to0$. If
    their dimensions agree but they are not gauge-phase equivalent, the
    equal-dimension overlap-decay theorem gives the same conclusion. In the
    remaining case, $B^i=\zeta X A^iX^{-1}$. By
    Theorem~\ref{thm:cpsv_normal_mpv_phase_gauge}, $|\zeta|=1$. Hence
    \[
        O_{AB}(N)=\overline{\zeta}^{\,N}O_{AA}(N),
        \qquad |O_{AB}(N)|=|O_{AA}(N)|\longrightarrow1.
    \]
\end{proof}

\begin{theorem}[Exact phase alternative for normal tensors]
    \label{thm:cpsv_normal_mpv_phase_alternative}
    \lean{MPSTensor.IsNormalTensor.mpv_phase_alternative}
    \leanok
    \uses{def:nt_cpsv, def:mpv_overlap, def:mpv_state}
    Let $A$ and $B$ be normal tensors of positive, possibly different, bond
    dimensions. Then either $O_{AB}(N)\to0$, or there is a phase
    $\zeta\in\C$, with $|\zeta|=1$, such that for every chain length $N$,
    \[
        \ket{V^{(N)}(B)}=\zeta^N\ket{V^{(N)}(A)}.
    \]
    This is \cite[Corollary~A.3]{Cirac2016MPDO_arXiv}.
\end{theorem}

\begin{proof}\leanok
    \uses{thm:cpsv_normal_overlap_dichotomy,
        thm:gauge_phase_mpv_scaling}
    Apply Theorem~\ref{thm:cpsv_normal_overlap_dichotomy}. In its second
    alternative, gauge-phase covariance gives
    $V^{(N)}(B)_\sigma=\zeta^N V^{(N)}(A)_\sigma$ for every configuration
    $\sigma$, and hence the asserted equality of vectors.
\end{proof}

\begin{theorem}[Unit overlap from eventual proportionality]
    \label{thm:eventual_proportional_normalized_overlap_norm_one}
    \lean{MPSTensor.mpvOverlap_norm_tendsto_one_of_eventually_proportionalMPV₂}
    \leanok
    \uses{def:mpv_overlap}
    Let $A$ and $B$ be tensors, possibly with different bond dimensions, whose
    self-overlaps tend to one. Suppose that, for all sufficiently large $N$,
    there is $c_N\in\C$ such that
    \[
        \ket{V^{(N)}(A)}=c_N\ket{V^{(N)}(B)}.
    \]
    Then $|O_{AB}(N)|\to1$.
\end{theorem}

\begin{proof}\leanok
    The proportionality relation gives
    \[
        O_{AB}(N)=c_NO_{BB}(N),
        \qquad O_{AA}(N)=|c_N|^2O_{BB}(N).
    \]
    Hence $|c_N|^2\to1$ and therefore $|c_N|\to1$.  The first identity now
    gives $|O_{AB}(N)|\to1$.
\end{proof}

\begin{theorem}[Geometric proportionality phase for normal tensors]
    \label{thm:normal_proportional_mpv_geometric_phase}
    \lean{MPSTensor.EventuallyNonzeroProportionalMPV₂.%
        exists_unit_phase_power_of_isNormalTensor}
    \leanok
    \uses{def:nt_cpsv, def:mpv_state,
        def:eventually_nonzero_proportional_mpv}
    Let $A$ and $B$ be normal tensors of positive, possibly different, bond
    dimensions. Suppose that, for all sufficiently large $N$, their matrix
    product vectors are proportional by a nonzero scalar. Then there is
    $a\in\C$, with $|a|=1$, such that, for every positive $N$,
    \[
        \ket{V^{(N)}(A)}=a^N\ket{V^{(N)}(B)}.
    \]
    This is the one-normal-block consequence of
    \cite[Corollary~A.3 and the proof of Theorem~2.10]{Cirac2016MPDO_arXiv}.
    % Maintainer note: this determines only the unit-phase part for normalized
    % normal blocks. The distinct positive-rescaling problem is recorded in
    % docs/paper-gaps/cpsv16_commuting_form_to_sal.tex.
\end{theorem}

\begin{proof}\leanok
    \uses{thm:eventual_proportional_normalized_overlap_norm_one,
        thm:cpsv_normal_mpv_phase_alternative,
        thm:cpsv_normal_mpv_phase_gauge}
    Normality gives $O_{AA}(N)\to1$ and $O_{BB}(N)\to1$, so
    Theorem~\ref{thm:eventual_proportional_normalized_overlap_norm_one} gives
    $|O_{AB}(N)|\to1$. The first alternative in
    Theorem~\ref{thm:cpsv_normal_mpv_phase_alternative} gives
    $O_{AB}(N)\to0$, hence $|O_{AB}(N)|\to0$, a contradiction. In the
    remaining alternative,
    $\ket{V^{(N)}(B)}=\zeta^N\ket{V^{(N)}(A)}$ with $|\zeta|=1$; take
    $a=\zeta^{-1}$.
\end{proof}

\begin{corollary}[Positive-length geometric proportionality phase]
    \label{cor:normal_positive_proportional_mpv_geometric_phase}
    \lean{MPSTensor.NonzeroProportionalMPV₂.%
        exists_unit_phase_power_of_isNormalTensor}
    \leanok
    \uses{def:nt_cpsv, def:mpv_state,
        def:nonzero_proportional_mpv}
    Let $A$ and $B$ be normal tensors of positive, possibly different, bond
    dimensions. Suppose that, for every positive $N$, their matrix product
    vectors are proportional by a nonzero scalar. Then there is $a\in\C$,
    with $|a|=1$, such that, for every positive $N$,
    \[
        \ket{V^{(N)}(A)}=a^N\ket{V^{(N)}(B)}.
    \]
    This is the positive-length one-normal-block consequence of
    \cite[Corollary~A.3 and Theorem~2.10]{Cirac2016MPDO_arXiv}.
\end{corollary}

\begin{proof}\leanok
    \uses{lem:nonzero_proportional_mpv_basic,
        thm:normal_proportional_mpv_geometric_phase}
    Positive-length nonzero proportionality implies eventual nonzero
    proportionality. Apply
    Theorem~\ref{thm:normal_proportional_mpv_geometric_phase}.
\end{proof}

\begin{theorem}[Separated normal tensors are eventually independent]
    \label{thm:cpsv_normal_separated_eventually_li}
    \lean{MPSTensor.exists_eventually_linearIndependent_of_normalTensor_blocks_not_gaugePhaseEquiv}
    \leanok
    \uses{def:nt_cpsv, def:blocks_not_gauge_phase_equiv}
    Let $(A_j)_{j=1}^g$ be a finite family of positive-bond-dimension normal
    tensors, no two of which are gauge-phase equivalent after identifying
    equal bond dimensions. Then the matrix product vectors
    $\{\ket{V^{(N)}(A_j)}\}_{j=1}^g$ are linearly independent for all
    sufficiently large $N$, as in
    \cite[Appendix~A]{Cirac2016MPDO_arXiv}.
\end{theorem}

\begin{proof}\leanok
    \uses{thm:cpsv_normal_tp_gauge,
        lem:bnt_finite_index_overlap_orthonormal_li}
    Put every tensor in the pure trace-preserving gauge of
    Theorem~\ref{thm:cpsv_normal_tp_gauge}. Distinct gauge-phase classes have
    overlaps tending to zero, while each self-overlap tends to one. The finite
    Gram matrix therefore tends to the identity and is invertible for all
    sufficiently large lengths. Gauge invariance of matrix product vectors
    gives the conclusion for the original tensors.
\end{proof}

\begin{theorem}[Simultaneous one-letter span after blocking a BNT]
    \label{thm:cpsv_bnt_blocked_simultaneous_span}
    \lean{MPSTensor.IsCPSVBasisOfNormalTensors.%
        exists_positive_wordTupleSpanTop}
    \lean{MPSTensor.wordTupleSpanTop_blockTensor_one}
    \lean{MPSTensor.IsCPSVBasisOfNormalTensors.%
        exists_blocked_wordTupleSpanTop_one}
    \leanok
    \uses{def:bnt_cpsv, def:blocked_tensor,
        def:blockwise_product_word_span}
    Let $(A_j)_{j=1}^g$ be a basis of normal tensors, with bond dimensions
    $D_j$. There is a positive integer $L$ such that the one-letter
    evaluations of the blocked tensors span the full product algebra:
    \[
        \spn_{\C}\left\{
        \bigl((A_1^{[L]})^i,\ldots,(A_g^{[L]})^i\bigr)
        : 0\leq i<d^L
        \right\}
        = \prod_{j=1}^g \MN{D_j}.
    \]
    This is the existence form of simultaneous block injectivity associated
    with \cite[lines~317--345]{Cirac2016MPDO_arXiv}; no numerical bound on
    $L$ is asserted here.
\end{theorem}

\begin{proof}\leanok
    \uses{thm:cpsv_normal_tp_gauge,
        thm:normal_tp_primitive_irred,
        thm:common_exact_block_injectivity_normal_left_canonical,
        thm:wordspan_blk_inj_succ,
        thm:peripheral_primitive_irreducible_overlap_one,
        lem:bnt_direct_sum_injective_prefix_word_span,
        lem:block_separating_equations_from_pairwise_separation,
        lem:product_word_span_from_bnt_block_separation,
        lem:gauge_invariance_block_injective}
    Eventual linear independence implies $D_j>0$ for every $j$ and excludes
    gauge-phase equivalence between two distinct blocks. Put every block in a
    trace-preserving gauge. Normality gives a common positive length $p$ at
    which every gauged block is $p$-block injective; this injectivity persists
    at all larger lengths. Pairwise direct-sum separation at lengths $p$,
    $p+1$, and $3(p+1)$ gives, for each ordered pair of distinct blocks, a
    length-$3(p+1)$ polynomial equal to the identity on the first block and
    zero on the second. Multiplying the $g-1$ pairwise polynomials gives a
    selector for each block. Combining these selectors with the common
    length-$p$ word spans yields the full product-algebra span at
    \[
        L=p+(g-1)3(p+1).
    \]
    Conjugation by the Perron gauges preserves this span. Finally, a one-letter
    evaluation of $A_j^{[L]}$ is precisely a length-$L$ word evaluation of
    $A_j$.
\end{proof}

Let $A$ be a tensor in canonical form, and let $A_1,\ldots,A_g$ be a basis of
normal tensors of $A$. For suitable coefficients $\mu_{j,q}$ and invertible
matrices $X_{j,q}$, write
$M_j=\diag(\mu_{j,1},\ldots,\mu_{j,r_j})$ and
$X=\bigoplus_{j=1}^{g}\bigoplus_{q=1}^{r_j}X_{j,q}$. Then
% Source: CPSV16, Section 2.3, labels eq:II_ABasicTensors and eq:II_X,
% source lines 283--301; rendered Equations (20a)--(21).
% Source-faithfulness verdict: the display labelled Eq19 in the source replaces
% direct sums by tensor products and is explicitly declared nonliteral at
% source lines 304--313; it is rendered as Equation (23).
% Index/type verdict: A^i_{alpha,beta} has physical index i and virtual indices
% alpha,beta; j and q select direct-sum summands and are not contraction legs.
\[
    \begin{array}{rcl}
        A^i
        &=&X\left[\bigoplus_{j=1}^{g}
            \left(M_j\otimes A_j^i\right)\right]X^{-1} \\
        &=&\bigoplus_{j=1}^{g}\bigoplus_{q=1}^{r_j}
            \mu_{j,q}X_{j,q}A_j^iX_{j,q}^{-1}.
    \end{array}
\]
This is \cite[Section~2.3, Equations~(20a)--(21)]{Cirac2016MPDO_arXiv}. It is
the decomposition used below to separate the normal blocks.

The eventual linear-independence clause is produced from the asymptotic
overlaps: diagonal overlaps tending to one and off-diagonal overlaps tending to
zero make the Gram matrix converge to the identity.

\begin{lemma}[Eventual linear independence from finite-index overlap orthonormality]
    \label{lem:bnt_finite_index_overlap_orthonormal_li}
    \lean{MPSTensor.eventually_linearIndependent_of_finite_overlap_tendsto_orthonormal}
    \leanok
    \uses{def:mpv_overlap}
    Let $I$ be a finite index set and let $(A_i)_{i \in I}$ be a finite
    family of tensors such that $O_{A_iA_i}(N) \to 1$ for each $i$, and
    $O_{A_iA_j}(N) \to 0$ whenever $i \neq j$. Then the MPV states
    $\{\ket{V^{(N)}(A_i)}\}_{i \in I}$ are linearly independent for all
    sufficiently large~$N$.
\end{lemma}

\begin{proof}\leanok
    The Gram matrix of the family $\{\ket{V^{(N)}(A_i)}\}_{i \in I}$
    converges entrywise to the identity matrix. Hence for all sufficiently
    large $N$ it is invertible, so the MPV states are linearly independent.
\end{proof}

\begin{lemma}[Eventual linear independence from overlap orthonormality]
    \label{lem:bnt_overlap_orthonormal_li}
    \lean{MPSTensor.eventually_linearIndependent_of_overlap_tendsto_orthonormal}
    \leanok
    \uses{def:mpv_overlap, lem:bnt_finite_index_overlap_orthonormal_li}
    Let $(A_j)_{j=1}^g$ be a finite family of tensors such that
    $O_{A_jA_j}(N) \to 1$ for each $j$ and $O_{A_iA_j}(N) \to 0$
    for $i \neq j$. Then the MPV states
    $\{\ket{V^{(N)}(A_j)}\}_{j=1}^g$ are linearly independent for all
    sufficiently large~$N$.
\end{lemma}

\begin{proof}\leanok
    Apply Lemma~\ref{lem:bnt_finite_index_overlap_orthonormal_li} with
    $I = \{1,\ldots,g\}$.
\end{proof}

The same conclusion is often more convenient with an explicit threshold rather
than the eventual quantifier.

\begin{theorem}[Existential BNT independence from overlap limits]%
    \label{thm:bnt_li_from_overlap_limits_exists}
    \lean{MPSTensor.exists_eventually_linearIndependent_of_overlap_tendsto_orthonormal}
    \leanok
    \uses{lem:bnt_overlap_orthonormal_li}
    If the self-overlaps of a finite block family tend to $1$ and the
    cross-overlaps of distinct blocks tend to $0$, then the MPV states are
    linearly independent for all sufficiently large system sizes.
\end{theorem}

\begin{proof}\leanok
    \uses{lem:bnt_overlap_orthonormal_li}
    This is Lemma~\ref{lem:bnt_overlap_orthonormal_li}, restated with an
    explicit threshold $N_0$.
\end{proof}

Eventual linear independence is also what converts an equality of two finite
linear combinations into equality of their coefficients; this is the mechanism
behind the coefficient identities of Chapter~\ref{ch:ft_proof}.

\begin{lemma}[Eventual coefficient extraction from linear independence]%
    \label{lem:eventual_coeff_eq_of_eventual_li}
    \lean{MPSTensor.coefficient_eventually_eq_of_eventually_linearIndependent}
    \leanok
    If a finite family of vectors is linearly independent for all sufficiently
    large $N$, and two finite linear combinations of that family are equal
    for all sufficiently large $N$, then the corresponding coefficients are
    equal for all sufficiently large $N$.
\end{lemma}

\begin{proof}\leanok
    Move the two linear combinations to one side. Linear independence forces
    each coefficient difference to vanish at every sufficiently large length.
\end{proof}

When two BNT families must be compared against each other, the same
Gram-matrix criterion applies to their disjoint union, provided the mixed
overlaps also decay. This is the form used in the exact sector matching of
Section~\ref{sec:sector_bnt_strong_match}.

\begin{lemma}[Eventual linear independence for two families]%
    \label{lem:bnt_two_family_overlap_orthonormal_li}
    \lean{MPSTensor.eventually_linearIndependent_of_two_family_overlap_tendsto_orthonormal}
    \leanok
    \uses{def:mpv_overlap, lem:bnt_finite_index_overlap_orthonormal_li}
    Let $(A_j)_{j=1}^{g_A}$ and $(B_k)_{k=1}^{g_B}$ be finite families of
    tensors. Suppose the overlaps inside each family are asymptotically
    orthonormal, and suppose every mixed overlap
    $O_{A_jB_k}(N)$ tends to~$0$. Then the combined MPV family
    \[
        \{\ket{V^{(N)}(A_j)}\}_{j=1}^{g_A}
        \cup
        \{\ket{V^{(N)}(B_k)}\}_{k=1}^{g_B}
    \]
    is linearly independent for all sufficiently large~$N$.
\end{lemma}

\begin{proof}\leanok
    Apply the finite-index Gram-matrix criterion to the disjoint union of the two
    index sets. The within-family hypotheses give the diagonal and
    off-diagonal limits inside each block, and the mixed overlap hypothesis gives
    the off-diagonal limits between the two blocks.
\end{proof}

We now isolate the separation condition that supplies the off-diagonal decay
and combine it with the canonical-form hypotheses.

\begin{definition}[No gauge-phase-equivalent distinct blocks]%
    \label{def:blocks_not_gauge_phase_equiv}
    \lean{MPSTensor.BlocksNotGaugePhaseEquiv}
    \leanok
    \uses{def:gauge_phase_equiv}
    A finite block family has no gauge-phase-equivalent distinct blocks when,
    for any two distinct indices with the same bond dimension, the corresponding
    blocks are not gauge-phase equivalent.
\end{definition}

\begin{definition}[Normal canonical form with BNT separation]\label{def:normal_cf_bnt}
    \lean{MPSTensor.IsNormalCanonicalFormBNT}\leanok
    \uses{def:is_normal_canonical_form, def:blocks_not_gauge_phase_equiv}
    A family is in \emph{normal canonical form with BNT separation} if it
    satisfies Definition~\ref{def:is_normal_canonical_form} and distinct blocks
    of the same bond dimension are not gauge-phase equivalent. The weight moduli
    are non-increasing but need not be strictly decreasing, so equal-modulus
    blocks are allowed; the grouping into a BNT is governed by gauge-phase
    equivalence, not by distinctness of moduli. Repeated equal-modulus copies of
    one basis tensor are not separate blocks here; they appear as the weights
    $\mu_{j,q}$ in $c_N^{(j)} = \sum_q \mu_{j,q}^N$ in the sector decomposition
    below.
    % Maintainer note: this is not the source two-layer BNT expansion
    %   A^i = \bigoplus_{j=1}^g \bigoplus_{q=1}^{r_j} \mu_{j,q} X_{j,q} A_j^i X_{j,q}^{-1},
    % where repeated copies stay visible through their coefficients,
    % multiplicities, and gauges. That expansion is recorded on the sector
    % decomposition (Definition~\ref{def:sector_bnt_cf}).
\end{definition}

The off-diagonal decay needed for the BNT criterion follows from the separation
condition together with the overlap-decay theorems of the preceding chapter.

\begin{theorem}[Cross-overlap decay from explicit BNT hypotheses]
    \label{thm:cross_overlap_zero_split}
    \lean{MPSTensor.cross_overlap_tendsto_zero_of_separated_bnt_data}
    \leanok
    \uses{def:mpv_overlap, def:gauge_phase_equiv}
    Let $(A_j)_{j=1}^r$ be a family in which each $A_j$ is injective and
    left-canonical, $\sum_i (A_j^i)^\dagger A_j^i = \Id$, and distinct
    equal-dimension blocks are not gauge-phase equivalent. Then
    $O_{A_jA_k}(N) \to 0$ for every pair $j \neq k$.
\end{theorem}

\begin{proof}\leanok\uses{thm:overlap_decay, thm:overlap_decay_rect}
    If the bond dimensions differ, Theorem~\ref{thm:overlap_decay_rect}
    gives the decay. If the bond dimensions agree, the non-equivalence
    hypothesis excludes gauge-phase equivalence, so
    Theorem~\ref{thm:overlap_decay} gives the decay.
\end{proof}

\begin{lemma}[BNT from explicit blockwise hypotheses]
    \label{lem:cf_bnt_is_bnt_split}
    \lean{MPSTensor.isBNT_of_separated_bnt_data}
    \leanok
    \uses{def:bnt, def:mpv_overlap}
    Let $(\mu_k, A_k)_{k=1}^r$ be a weighted block family. Assume each $A_k$ is
    injective and left-canonical, $\sum_i (A_k^i)^\dagger A_k^i = \Id$, that
    $O_{A_kA_k}(N) \to 1$ for each $k$, and that distinct equal-dimension blocks
    are not gauge-phase equivalent. Then the blocks form a basis of normal
    tensors for the block-diagonal tensor $\bigoplus_k \mu_k A_k$.
\end{lemma}

\begin{proof}\leanok
    \uses{thm:mpv_decomposition, lem:bnt_overlap_orthonormal_li, thm:cross_overlap_zero_split}
    Injective blocks are normal. The block-diagonal decomposition theorem gives
    the MPV expansion of $\bigoplus_k \mu_k A_k$ in terms of the block MPVs.
    The self-overlap limits are part of the hypotheses, and
    Theorem~\ref{thm:cross_overlap_zero_split} supplies the off-diagonal
    decay. Therefore Lemma~\ref{lem:bnt_overlap_orthonormal_li} gives the
    eventual linear independence required in Definition~\ref{def:bnt}.
\end{proof}

The normal-canonical formulation encodes normality spectrally (primitive
transfer map), the BNT definition algebraically (eventual block injectivity).

\begin{remark}[Comparing the two normality conditions]%
    \label{rem:normalcf_bnt_gap}
    The spectral normality of Definition~\ref{def:is_normal_canonical_form} and
    the algebraic normality of Definition~\ref{def:normal} agree: the Wielandt
    theory of Chapter~\ref{ch:wielandt} supplies the implication from
    primitivity to eventual block injectivity used here.
\end{remark}

\begin{lemma}[Restricted normal-CF-BNT hypotheses give a BNT]
    \label{lem:normal_cf_bnt_is_bnt}
    \lean{MPSTensor.IsNormalCanonicalFormBNT.isBNT}
    \leanok
    \uses{def:normal_cf_bnt, def:bnt}
    If a family is in normal canonical form with BNT separation and each block
    is also normal in the algebraic sense, then the blocks form a basis of
    normal tensors for the associated block-diagonal tensor. The blocks here are
    distinct representatives.
    % Scope restriction (ft_one_copy_scope_restriction): this form carries no
    % repeated-copy multiplicities; the missing multiplicity argument is
    % recorded in docs/paper-gaps/ft_one_copy_scope_restriction.tex.
\end{lemma}

\begin{proof}\leanok\uses{rem:normalcf_bnt_gap}
    The normal-canonical hypotheses already give the block-diagonal MPV
    expansion, the self-overlap normalization, and the eventual linear
    independence coming from irreducible trace-preserving overlap decay.
    Supplying algebraic normality for each block adds the remaining clause in
    Definition~\ref{def:bnt}.
\end{proof}
